\section{基于运行生产资料的汽轮机系统故障知识图谱构建}
\label{sec:knowledge_graph_v6}

\subsection{引言}

燃煤机组汽轮机系统由多级汽缸、再热系统、抽汽回热系统以及凝汽器冷端等多个热力设备组成，各设备通过蒸汽、凝结水和给水持续进行物质与能量传递。局部设备发生通流、汽封、阀门或换热性能异常后，其影响通常不会局限于单一监测参数，而会沿实际热力连接向相邻设备和参数传播。第二章建立的健康状态模型能够在给定运行边界下获得设备状态参数的健康参考值，并由实测值与健康预测值之间的偏差表征状态变化程度；但单独依赖参数偏差，仍难以刻画参数所属设备、参数之间的物理联系以及多参数异常与故障机理之间的关联。

电站长期设计、运行和维护过程中积累了大量设备结构、DCS测点、系统接口、故障分析和检修经验等生产资料。这些知识既存在于点表、设备清单等结构化数据中，也大量分布于运行说明、技术报告和公开文献等非结构化文本中。知识图谱（Knowledge Graph，KG）能够以实体和关系的形式组织多源异构知识，并利用图结构显式描述设备、参数、异常和故障之间的语义联系\cite{Hogan2021KnowledgeGraphs}。因此，将汽轮机运行生产资料转化为领域知识图谱，可为后续图卷积故障诊断提供区别于纯统计相关性的结构先验。

本章面向燃煤机组汽轮机系统构建基于运行生产资料的故障知识图谱。首先分析汽轮机系统参数关联与典型故障特性，明确图谱需要描述的知识对象；随后整理机组DCS点表、设备输入输出关系和系统边界资料，并利用公开论文、标准及运行事件资料补充故障机理知识；在此基础上定义知识图谱节点与关系类型，采用大语言模型生成候选三元组，通过领域规则、证据来源和Embedding语义对齐完成知识校核与实体标准化，最终形成可与实际DCS状态量建立映射的汽轮机系统故障知识图谱。

\subsection{燃煤机组汽轮机系统参数及故障特性}
\label{subsec:turbine_fault_characteristics_v6}

汽轮机主通流可概括为“主蒸汽--VHP--一次再热--HP--二次再热--IP--LP--凝汽器”。主蒸汽压力、温度和流量构成VHP入口热力边界，VHP排汽进入一次再热系统，一次热再蒸汽形成HP入口条件，HP排汽经二次再热后进入IP，最终蒸汽经LP膨胀后排入凝汽器。除主通流外，各级汽缸还通过抽汽支路向回热系统提供热源，使汽轮机本体与高压加热器、低压加热器、除氧器和给水系统形成进一步耦合。根据本文机组设备IO与DCS资料，VHP一级抽汽与1号高压加热器相连，HP二级抽汽与2号高压加热器相连，IP三级和四级抽汽分别与3号、4号高压加热器相连，LP六级至八级抽汽分别与相应低压加热器支路连接。对于现有资料中尚未形成一致描述的局部接口关系，不将其作为确定性拓扑写入后续参数图。

汽轮机热力故障具有显著的多参数耦合特征。固体颗粒侵蚀会改变喷嘴或叶片有效通流面积，并进一步影响关键压力分布\cite{Kubiak2002TurbineDegradation}；叶片表面沉积可造成有效通流面积减小并引起效率和出力下降\cite{Kubiak2002ScaleDeposit}；级间或轴端汽封缺陷会导致蒸汽泄漏，低压侧密封异常还可能伴随空气漏入和凝汽器真空恶化\cite{ZhangHu2021FaultPropagationGNN}；HP--IP中间汽封泄漏与IP效率、一次热再温度及IP排汽温度等状态量具有诊断关联\cite{Xu2011HPIPLeakage}；凝汽器水侧污垢会降低换热能力并影响背压及低压缸可利用焓降\cite{Li2016CondenserFouling,Medica2018CondenserCondition}；给水加热器内部泄漏则会破坏正常回热过程，在抽汽侧和给水侧形成协同状态变化\cite{Heo2012FWHLeakage,Barszcz2011FeedwaterHeater}。

上述故障具有两方面共同特征：其一，同一故障在多个参数上的响应受设备归属和热力传播路径约束；其二，不同故障可能共享部分压力、温度或流量参数，仅依据局部数值变化容易产生类别混淆。因此，故障知识图谱不仅需要记录“故障与参数相关”的静态事实，还需要同时表达设备连接、参数归属、故障发生位置、异常表现及处置关系，从而为多参数故障诊断提供结构化知识基础。

\subsection{基于运行生产资料的知识图谱构建方法}
\label{subsec:kg_construction_v6}

\subsubsection{汽轮机热力过程的运行生产资料}
\label{subsubsec:operation_materials_v6}

运行生产资料是领域故障知识图谱的主要知识来源。本文能够直接获得的机组资料包括DCS测点表、设备输入输出关系、系统边界核对结果以及状态监测建模资料。该类资料能够确定设备名称、参数点号、参数归属、主通流方向、抽汽支路及相邻设备接口，为知识图谱中的设备和参数实体提供机组侧事实依据。

现场结构资料通常难以完整覆盖典型故障的原因、演化过程和异常表现。为补充故障机理知识，进一步引入同行评议论文、标准、设备运行维护资料和公开运行事件报告。跨机组资料主要用于提取具有通用物理意义的“故障--异常--参数”关系，而不直接将其他机组某一具体测点的变化方向套用于本文机组。例如，公开研究能够支持“汽封缺陷导致蒸汽泄漏”这一通用机理，但当文献未对本文机组某一DCS温度测点给出直接证据时，仅将其作为待解释的局部响应，而不写成确定性点级关系。主要资料来源及用途见表~\ref{tab:kg_sources_v6}。

\begin{table}[htbp]
    \centering
    \caption{汽轮机故障知识图谱的主要知识来源}
    \label{tab:kg_sources_v6}
    \begin{tabular}{p{3.8cm}p{5.2cm}p{5.2cm}}
        \toprule
        资料类型 & 主要内容 & 图谱中的主要用途 \\
        \midrule
        DCS点表与设备IO资料 & 测点名称、点号、设备归属、入口和出口边界 & 构建设备、参数实体及机组直接拓扑 \\
        系统边界与状态建模资料 & 主通流、抽汽、再热及设备边界关系 & 校核设备接口和DCS映射 \\
        同行评议论文 & 汽封泄漏、通流侵蚀、沉积、冷端及回热故障机理 & 提取故障、异常和参数之间的通用关系 \\
        标准、维护资料及公开运行事件 & 故障现象、运行影响和处置经验 & 补充并交叉核验故障关系 \\
        \bottomrule
    \end{tabular}
\end{table}

在知识入图过程中，对机组结构事实与外部故障机理进行分层管理。由DCS、设备IO和系统拓扑直接确认的关系用于描述本文机组设备结构；由论文、标准或公开运行事件直接支持的关系用于描述通用故障机理；仅能依据热力规律推导而缺少直接证据的点级响应不作为主知识图谱中的确定关系。该处理既保证设备拓扑与本文机组一致，也避免将跨机组经验过度外推到具体DCS测点。

\subsubsection{节点和关系的构建与提取}
\label{subsubsec:entity_relation_v6}

运行生产资料中的设备、参数和故障信息通常以自然语言、工程简称或DCS编码等不同形式出现，同一对象在不同资料中可能存在多种表达。为使多源知识能够在统一图结构中表示，首先定义知识图谱实体类型和关系类型。设汽轮机系统故障知识图谱为
\begin{equation}
G_{KG}=(V,E,\mathcal R),
\end{equation}
式中：$V$为实体集合；$E$为实体间关系边集合；$\mathcal R$为关系类型集合。根据汽轮机系统故障诊断需求，将实体划分为设备、参数、异常、故障和处置五类：
\begin{equation}
V=V_{dev}\cup V_{par}\cup V_{abn}\cup V_{fault}\cup V_{act}.
\end{equation}
设备实体表示VHP、HP、IP、LP、再热系统、抽汽支路、加热器和凝汽器等物理对象；参数实体表示可由DCS直接获取的压力、温度、流量和阀位等运行量；异常实体描述蒸汽泄漏、真空下降、通流面积变化和效率降低等状态现象；故障实体表示具有独立工程含义的故障类型；处置实体表示检查、清洗和检修等维护措施。

根据上述实体类型定义设备归属、介质流向、设备接口、发生于、导致、表现为、影响、关联监测和处置为等关系，主要关系类型如表~\ref{tab:kg_relation_types_v6}所示。

\begin{table}[htbp]
    \centering
    \caption{汽轮机故障知识图谱主要关系类型}
    \label{tab:kg_relation_types_v6}
    \begin{tabular}{lll}
        \toprule
        关系类型 & 头实体类型 & 尾实体类型 \\
        \midrule
        归属 & 参数 & 设备 \\
        介质流向/设备接口 & 设备 & 设备 \\
        发生于 & 故障 & 设备 \\
        导致/表现为 & 故障或异常 & 异常 \\
        影响/关联监测 & 故障或异常 & 参数 \\
        处置为 & 故障 & 处置 \\
        \bottomrule
    \end{tabular}
\end{table}

在候选知识抽取阶段，首先利用大语言模型对设备资料和故障文本进行语义解析，并按照预定义实体与关系Schema输出候选三元组。该过程能够快速识别文本中的设备、参数、异常、故障及处置对象，但语言模型输出仅作为候选知识，仍需结合领域规则和证据来源进行校核。领域规则主要包括实体类型约束、关系方向约束、设备兼容性约束及参数可观测性约束。例如，“压力升高”“阀位异常”等带有状态方向的表述应归入异常实体，而不能直接作为实时数值参数节点；不存在于本文机组结构中的设备连接，也不能仅依据通用汽轮机结构建立为机组确定关系。

在实体标准化阶段，采用Embedding语义对齐处理设备简称、参数别名和不同文本写法\cite{Reimers2019SentenceBERT}。对于候选实体$v_i$和标准实体$v_j$，综合语义向量相似度和字符串相似度计算匹配程度：
\begin{equation}
S_{ij}=\lambda S_{emb}(v_i,v_j)+(1-\lambda)S_{str}(v_i,v_j),
\end{equation}
式中：$S_{emb}$为Embedding语义相似度；$S_{str}$为字符串相似度；$\lambda$为两类相似度的权重系数。当综合相似度达到设定阈值且实体类型一致时，将候选实体归并至标准实体；否则保留为待核实体。

为综合衡量候选关系的可靠性，对语言模型抽取结果、领域规则、实体标准化结果和来源证据进行联合评价。对于第$e$条候选关系，其置信度可表示为
\begin{equation}
C_e=\alpha C_{llm}+\beta C_{rule}+\gamma C_{emb}+\delta C_{evi},
\end{equation}
式中：$C_{llm}$为语言模型候选抽取置信度；$C_{rule}$为领域规则一致性得分；$C_{emb}$为实体语义对齐置信度；$C_{evi}$为来源证据强度；$\alpha$、$\beta$、$\gamma$和$\delta$为相应权重。机组设备连接关系最终以DCS和设备IO资料为优先依据，故障机理关系则需保留可追溯的论文、标准或运行事件来源。经过候选抽取、规则校核和实体标准化后，分散于多源资料中的知识被统一表示为“头实体--关系--尾实体”三元组，并进一步组成汽轮机系统领域知识图谱。

\subsection{知识图谱构建结果}
\label{subsec:kg_result_v6}

按照上述流程，形成了覆盖汽轮机主通流、再热、抽汽回热和凝汽冷端等对象的故障知识库。当前版本共包含127个标准实体和123条主关系，其中81条关系来自本文机组设备结构与DCS资料，40条关系由公开论文、标准或运行事件中的直接故障机理支撑；另有2条机组资料之间存在描述差异的候选关系未纳入主知识图谱。由此，最终图谱同时包含机组设备结构知识和具有来源依据的通用故障机理知识。

在设备拓扑方面，图谱保存“主蒸汽--VHP--一次再热--HP--二次再热--IP--LP--凝汽器”的主通流关系，并包含VHP一级抽汽至1号高压加热器、HP二级抽汽至2号高压加热器、IP三级和四级抽汽至3号和4号高压加热器，以及LP相关抽汽至低压加热器支路等机组关系。在故障语义方面，图谱进一步连接汽封缺陷、蒸汽泄漏、通流侵蚀、叶片沉积、凝汽器污垢和回热设备泄漏等故障与相应异常状态、关联参数及处置知识。

系统级知识图谱用于表达完整汽轮机故障领域知识，而图卷积模型需要加载可实时获取的数值特征。因此，在知识图谱构建完成后，进一步建立标准参数实体与实际DCS测点之间的映射。第五章VHP代表性算例从系统知识图谱中抽取VHP本体、一级抽汽支路和一次再热接口对应的状态参数节点，形成设备任务参数子图，并将第二章健康状态模型产生的动态残差信息加载至对应参数节点，由此实现静态领域知识与动态运行数据之间的连接。

\subsection{本章小结}

本章面向燃煤机组汽轮机系统构建了基于运行生产资料的故障知识图谱。首先分析汽轮机主通流、再热、抽汽回热和凝汽冷端之间的参数关联以及典型热力故障的多参数传播特征；随后整理机组DCS点表、设备IO、系统边界资料，并利用公开论文、标准和运行事件补充故障机理知识；在此基础上定义设备、参数、异常、故障和处置等实体及相应关系类型，采用大语言模型生成候选三元组，并结合领域规则、证据来源和Embedding语义对齐完成知识校核与实体标准化。最终形成能够与实际DCS参数建立映射的汽轮机系统故障知识图谱，为下一章基于知识图谱和图卷积网络的故障诊断模型提供静态结构信息。